\documentclass[UTF8, 12pt, fontset=fandol, oneside]{ctexart}
\usepackage{researchnotes}

\title{伯努利不等式}
\author{}
\date{}

\begin{document}

\maketitle

伯努利不等式是初等不等式中最常用的工具之一. 它说明：当一个数在 $1$ 附近作相同方向的乘法累积时，所得结果至少会保留一次项所给出的线性变化. 这个不等式形式简单，但在估计幂、证明极限、处理增长速度和建立更复杂的不等式时都十分有用.

\section*{一、标准形式}

\noindent\textbf{定理 1 (伯努利不等式)}\quad 设 $x\ge -1$，$n$ 为正整数，则
\[
  (1+x)^n\ge 1+nx.
\]
当 $n=1$ 时，对所有 $x\ge -1$ 都取等号；当 $n>1$ 时，当且仅当 $x=0$ 取等号.

\noindent\textbf{说明}\quad 条件 $x\ge -1$ 的作用是保证 $1+x\ge0$. 在归纳证明中，我们需要把不等式两边同乘以 $1+x$，此时只有 $1+x\ge0$ 才不会改变不等号方向. 另外，从实际意义看，$1+x$ 往往表示某种倍率，因此自然应是非负的.

\section*{二、归纳法证明}

\noindent\textbf{证明}\quad 当 $n=1$ 时，
\[
  (1+x)^1=1+x=1+1\cdot x,
\]
结论成立.

假设对某个正整数 $n$，已有
\[
  (1+x)^n\ge 1+nx.
\]
由于 $x\ge -1$，有 $1+x\ge0$，所以
\[
\begin{aligned}
  (1+x)^{n+1}
  &= (1+x)^n(1+x)\\
  &\ge (1+nx)(1+x)\\
  &= 1+(n+1)x+nx^2\\
  &\ge 1+(n+1)x.
\end{aligned}
\]
因此，若结论对 $n$ 成立，则对 $n+1$ 也成立. 由数学归纳法，伯努利不等式对一切正整数 $n$ 成立.

当 $n>1$ 且 $x\ne0$ 时，上式最后一行中的 $nx^2>0$，并且从归纳过程可以看出不等式严格成立，故
\[
  (1+x)^n>1+nx.
\]
这也给出了等号条件.

\section*{三、二项式角度}

当 $x\ge0$ 时，伯努利不等式还可以直接由二项式定理看出：
\[
\begin{aligned}
  (1+x)^n
  &=1+nx+\binom n2x^2+\binom n3x^3+\cdots+x^n\\
  &\ge 1+nx.
\end{aligned}
\]
这个证明十分直观：展开式中除 $1+nx$ 外，其余各项都是非负项.

不过，这个证明只适合 $x\ge0$ 的情形. 当 $-1\le x<0$ 时，二项式展开中的高次项会正负交替，不能简单地说余项非负. 因此，完整的标准证明通常采用归纳法，或者采用下面的凸性方法.

\section*{四、凸性证明}

令 $y=1+x$，则条件 $x\ge -1$ 等价于 $y\ge0$，而要证明的不等式变为
\[
  y^n\ge 1+n(y-1).
\]
设
\[
  f(y)=y^n,\qquad y\ge0.
\]
当 $n\ge2$ 时，
\[
  f''(y)=n(n-1)y^{n-2}\ge0,
\]
所以 $f$ 在 $[0,+\infty)$ 上是凸函数. 凸函数的图像位于任一点切线的上方，取切点 $y=1$，便有
\[
  f(y)\ge f(1)+f'(1)(y-1)=1+n(y-1).
\]
于是
\[
  (1+x)^n\ge1+nx.
\]
当 $n=1$ 时结论显然成立. 这说明伯努利不等式可以看作幂函数 $y^n$ 在 $y=1$ 处的一阶切线估计.

\section*{五、实指数推广}

整数指数的伯努利不等式只是更一般事实的一个特例.

\noindent\textbf{定理 2 (实指数形式)}\quad 设 $x>-1$，$\alpha\in\mathbb R$.
\begin{enumerate}
  \item 若 $\alpha\ge1$ 或 $\alpha\le0$，则
  \[
    (1+x)^\alpha\ge1+\alpha x.
  \]
  \item 若 $0\le\alpha\le1$，则
  \[
    (1+x)^\alpha\le1+\alpha x.
  \]
\end{enumerate}
在非平凡情形 $\alpha\ne0,1$ 中，等号当且仅当 $x=0$.

\noindent\textbf{证明}\quad 设
\[
  f(t)=(1+t)^\alpha,\qquad t>-1.
\]
则
\[
  f'(t)=\alpha(1+t)^{\alpha-1},\qquad
  f''(t)=\alpha(\alpha-1)(1+t)^{\alpha-2}.
\]
若 $\alpha\ge1$ 或 $\alpha\le0$，则 $f''(t)\ge0$，所以 $f$ 是凸函数. 凸函数位于切线的上方，于是在 $t=0$ 处有
\[
  f(x)\ge f(0)+f'(0)x=1+\alpha x.
\]
若 $0\le\alpha\le1$，则 $f''(t)\le0$，所以 $f$ 是凹函数. 凹函数位于切线的下方，于是
\[
  f(x)\le f(0)+f'(0)x=1+\alpha x.
\]
证毕.

\noindent\textbf{例 1}\quad 取 $\alpha=\dfrac12$，当 $x>-1$ 时，由 $0<\alpha<1$ 得
\[
  \sqrt{1+x}\le 1+\frac{x}{2}.
\]
例如取 $x=a^2-1$，$a>0$，可得
\[
  a\le \frac{a^2+1}{2},
\]
即
\[
  a+\frac1a\ge2.
\]

\noindent\textbf{例 2}\quad 取 $\alpha=-1$，当 $x>-1$ 时，
\[
  \frac1{1+x}\ge1-x.
\]
更一般地，对正整数 $n$ 有
\[
  \frac1{(1+x)^n}\ge1-nx,\qquad x>-1.
\]

\section*{六、常见变形}

\noindent\textbf{变形 1}\quad 若 $a\ge0$，$n$ 为正整数，则
\[
  (1+a)^n\ge1+na.
\]
这是最常用的非负形式. 它说明固定比例的反复增长至少包含线性增长量.

\noindent\textbf{变形 2}\quad 若 $t\ge0$，$n$ 为正整数，则
\[
  \left(1+\frac{t}{n}\right)^n\ge1+t.
\]
这是把标准形式中的 $x$ 换成 $\dfrac tn$ 得到的. 这个形式常出现在复利、极限和指数函数的估计中.

\noindent\textbf{变形 3}\quad 若 $0<q<1$，则对任意正整数 $n$，
\[
  q^n\le \frac1{1+n\left(\frac1q-1\right)}.
\]
事实上，令 $x=\dfrac1q-1>0$，则 $\dfrac1q=1+x$. 由伯努利不等式，
\[
  q^{-n}=\left(1+x\right)^n\ge1+nx,
\]
两边取倒数即得结论. 这个估计常用来说明几何衰减会趋于 $0$.

\noindent\textbf{变形 4}\quad 若 $x\ge0$，$n$ 为正整数，则可保留二项式展开中的更多非负项：
\[
  (1+x)^n\ge1+nx+\binom n2x^2.
\]
一般地，对 $0\le k\le n$，
\[
  (1+x)^n\ge \sum_{j=0}^k\binom njx^j,\qquad x\ge0.
\]
这是伯努利不等式的加强版，但它依赖 $x\ge0$.

\section*{七、典型应用}

\noindent\textbf{例 3}\quad 证明对任意正整数 $n$，
\[
  2^n\ge n+1.
\]

\noindent\textbf{证明}\quad 在伯努利不等式中取 $x=1$，得
\[
  2^n=(1+1)^n\ge1+n.
\]

\noindent\textbf{例 4}\quad 设 $a>1$，证明数列 $a^n$ 随 $n$ 的增长至少线性趋于无穷.

\noindent\textbf{证明}\quad 记 $a=1+x$，其中 $x=a-1>0$. 由伯努利不等式，
\[
  a^n=(1+x)^n\ge1+nx.
\]
由于 $x>0$，当 $n\to+\infty$ 时，$1+nx\to+\infty$，所以
\[
  a^n\to+\infty.
\]

\noindent\textbf{例 5}\quad 设 $0<q<1$，证明 $q^n\to0$.

\noindent\textbf{证明}\quad 令
\[
  x=\frac1q-1>0.
\]
由伯努利不等式，
\[
  q^{-n}=\left(1+x\right)^n\ge1+nx.
\]
于是
\[
  0<q^n\le\frac1{1+nx}.
\]
令 $n\to+\infty$，右端趋于 $0$，由夹逼原理知
\[
  q^n\to0.
\]

\noindent\textbf{例 6}\quad 证明当 $n\ge2$ 时，
\[
  \left(1+\frac1n\right)^n\ge2.
\]

\noindent\textbf{证明}\quad 在伯努利不等式中取 $x=\dfrac1n$，得
\[
  \left(1+\frac1n\right)^n\ge1+n\cdot\frac1n=2.
\]

\noindent\textbf{例 7}\quad 设 $x\ge0$，证明
\[
  (1+x)^n-1\ge nx.
\]

\noindent\textbf{证明}\quad 这只是伯努利不等式的等价写法：
\[
  (1+x)^n\ge1+nx
  \quad\Longleftrightarrow\quad
  (1+x)^n-1\ge nx.
\]
它常用于估计幂函数增量，例如说明 $y=x^n$ 在正半轴上的增量至少含有一次主项.

\section*{八、与均值不等式的关系}

伯努利不等式也可以从算术平均数--几何平均数不等式中得到一个常见情形. 若 $x\ge -1$ 且 $1+nx\ge0$，则把 $n$ 个非负数
\[
  1+nx,\ 1,\ 1,\ \ldots,\ 1
\]
代入均值不等式，得
\[
  \frac{(1+nx)+(n-1)}{n}\ge
  \sqrt[n]{1+nx}.
\]
左边正好等于 $1+x$，所以
\[
  1+x\ge\sqrt[n]{1+nx}.
\]
两边取 $n$ 次幂，得到
\[
  (1+x)^n\ge1+nx.
\]
这种证明要求 $1+nx\ge0$，因此它只覆盖 $x\ge-\dfrac1n$ 的情形；而标准伯努利不等式适用于更大的范围 $x\ge-1$.

\section*{九、易错点}

\noindent\textbf{易错点 1}\quad 不要在 $-1\le x<0$ 时直接用二项式展开并声称余项全为非负. 此时 $x^3,x^5,\ldots$ 为负，余项的符号并不显然.

\noindent\textbf{易错点 2}\quad 实指数形式中，不等号方向取决于指数范围. 当 $0<\alpha<1$ 时，
\[
  (1+x)^\alpha\le1+\alpha x,
\]
方向与整数指数情形相反.

\noindent\textbf{易错点 3}\quad 若 $x<-1$，标准归纳证明失效，伯努利不等式也不再作为统一结论成立. 例如取 $x=-10$，$n=3$，则
\[
  (1+x)^n=(-9)^3=-729,\qquad 1+nx=1-30=-29,
\]
显然
\[
  -729\not\ge -29.
\]

\noindent\textbf{易错点 4}\quad 等号条件不要漏掉 $n=1$ 的特殊情形. 若 $n=1$，恒有等号；若 $n>1$，只有 $x=0$ 时取等号.

\section*{十、小结}

伯努利不等式的核心形式是
\[
  (1+x)^n\ge1+nx,\qquad x\ge-1,\quad n\in\mathbb N_+.
\]
它可以理解为幂函数在 $1$ 附近的一阶线性下估计. 对整数指数，归纳法给出最稳妥的证明；对实指数，凸性和凹性给出统一解释. 在应用中，它常用于把幂式估计为线性式，或把几何增长、几何衰减转化为容易处理的一次函数估计.

\end{document}
